\section{Eigenvalue Lens：从巨大 vocabulary matrix 压缩出 copying signature}

上一章已经把 head 折叠成 token-to-token matrices，但这只是把问题变得可定义，还没有让它变得可阅读：$M_{OV}$ 与 $M_{QK}$ 的边长等于 vocabulary size，逐项枚举既昂贵，也很难形成全局判断。本节引入 eigenvalue lens，核心问题是能否用少量 spectral signatures 判断一个 head 更像 copying、anti-copying，还是 different-token mapping；读图时还要始终记住，这种压缩会丢失具体 token、basis 与非线性上下文。

若矩阵把同一向量空间映射回自身，可以用 eigenvectors/eigenvalues 总结它在某些 directions 上是放大、反转还是旋转。

\lecturefigure{slide-26-eigenvector-introduction.jpg}{当输入输出属于同一空间时，eigenvectors 可压缩矩阵行为。}{Stanford Online 官方视频 00:33:58--00:34:21。}

\[
Mv=\lambda v.
\]

\begin{itemize}
  \item $v$：矩阵作用后方向保持不变的 mode；
  \item $\lambda$：该 mode 的缩放/相位变化；
  \item real positive $\lambda$：沿同一方向增强；
  \item real negative $\lambda$：反向抑制；complex $\lambda$：在二维子空间旋转并缩放。
\end{itemize}

\lecturefigure{slide-27-complex-eigenvalues.jpg}{非对称 token-to-token maps 可以拥有 complex eigenvalues。}{Stanford Online 官方视频 00:34:22--00:35:09。}

\subsection{OV spectrum 与 copying}

对 token-space OV map，正 eigenmodes 表示看到某组 token directions 后提高同类 output directions 的 logits，可解释为 copying tendency；负 modes 是 anti-copying；complex/imaginary modes 更偏向把一种 token direction 映射到另一种。

\lecturefigure{slide-28-ov-eigenvalue-meaning.jpg}{OV eigenvalue 图的教学解释：anti-copying、different-token mapping 与 copying。}{Stanford Online 官方视频 00:35:10--00:35:55。}

\begin{warningbox}{Spectrum 是 summary，不是完整语义字典}
Vocabulary embedding/unembedding 不一定正交，矩阵也可能 non-normal；eigenvectors 未必对应单个可命名 token。将正 eigenvalue 简化为 copying 是对 small attention-only circuits 的实用诊断，不是对任意矩阵的普遍语义定理。
\end{warningbox}

\subsection{一层与两层 heads 的 copying distribution}

Olah 展示一层模型中十个（共十二个）heads 的 OV spectrum 主要落在正实方向，说明模型大量使用“看到 token 后提升相同/相关 token”这一简单策略。两层模型中也有大量 copying heads，但多层 routing 允许它们被更精确的 pattern 触发。

\lecturefigure{slide-29-layer-one-copying-heads.jpg}{一层模型：10/12 attention heads 具有强 copying signature。}{Stanford Online 官方视频 00:35:56--00:36:05。}

\lecturefigure{slide-30-layer-two-copying-heads.jpg}{两层模型：第二层仍有多个强 copying heads。}{Stanford Online 官方视频 00:36:06--00:36:41。}

\lecturefigure{slide-31-head-eigenvalue-histogram.jpg}{把 OV spectral signature 压缩为全模型 head histogram。}{Stanford Online 官方视频 00:36:42--00:37:21。}

\teachervoice{讲者强调这是 bottom-up methodology：先从矩阵结构推导可读 summary，再说“这些 heads 多数在 copy”，而不是先假设所有 heads 都是某类 pattern 后去挑支持例子。}

\subsection{QK spectrum 与方法失效边界}

QK 的 positive mode 倾向匹配相同 token direction，negative mode 倾向避免相同 token，complex mode 倾向不同-token relation。它在多层模型中更有意义，因为 previous heads 可以把“前一个 token 是什么”等信息写入当前 residual，再由下一层 QK 做 same-feature matching。

\lecturefigure{slide-32-qk-eigenvalues-multilayer.jpg}{QK eigenvalues 在多层 chains 中更有用：前层先移动信息，后层再匹配。}{Stanford Online 官方视频 00:37:22--00:37:37。}

\lecturefigure{slide-33-eigenanalysis-breaks-at-scale.jpg}{Naive eigenvalue analysis 随模型增大和加入 MLP 而失效。}{Stanford Online 官方视频 00:37:38--00:38:35。}

\begin{warningbox}{MLP 让“一个矩阵解释一切”失效}
Attention-only model 在 fixed pattern 下可局部线性化；MLP 引入 activation-dependent nonlinear feature transforms，LayerNorm 也耦合坐标。大型模型仍可分析局部 Jacobian、features 与 paths，但不能直接把 small-model eigenvalue story 原样复制。
\end{warningbox}

\teachervoice{Olah 明确说 naive spectral method 在加入 MLP 后 “completely breaks down”。高质量讲义必须保留这句话，不能把 eigenvalue plots 包装成通用 LLM interpretability toolkit。}

\subsection{本章小结}

Eigenvalue lens 将巨大 token-to-token circuit 压缩成 copying/anti-copying 等趋势，并揭示一层模型强烈依赖广泛 copying。它依赖同空间映射和近似线性结构，面对 MLP 与大模型时只能作为有限线索。

